\documentclass[10pt]{beamer}
\usepackage[utf8]{inputenc}

\usetheme{Madrid}
\usecolortheme{default}
\usepackage{graphicx}
\usepackage{amsmath}
\usepackage{booktabs}
\usepackage{amsmath,amssymb,amsthm,microtype}
\usepackage{tikz}
\graphicspath{{../OverleafPaper/figures/}{OverleafPaper/figures/}{../forgetting/outputs/test/monte_carlo/figures/}{forgetting/outputs/test/monte_carlo/figures/}{../forgetting/outputs/test/first_2000/figures/individual_states/p0_0.5/}{forgetting/outputs/test/first_2000/figures/individual_states/p0_0.5/}{../forgetting/outputs/test/first_2000/figures/state_distributions/p0_0.5/}{forgetting/outputs/test/first_2000/figures/state_distributions/p0_0.5/}}

\title{Weekly report}
\author{Natalia Trigo Tomasevich}
\date{October 19, 2026}

\begin{document}

\frame{\titlepage}

\begin{frame}{This week}
    \begin{itemize}[<+->]
        \item Competition model
        \item Oblivious equilibrium 
        \item Simulations 
    \end{itemize}
\end{frame}

\begin{frame}{Elements of the model}
    \begin{itemize}
        \item Two symmetric providers, indexed by $j\in\{A,B\}$
        \item Two products, indexed by $i\in\{1,2\}$
        \item Success probabilities: $0<p_1<p_2<1$
        \item Costs: $c_1<c_2$
        \item Revenue: $R>c_2$
    \end{itemize}
\end{frame}

\begin{frame}{Customer's beliefs}
    For the competition model, the customer observes the state $x_A,x_B$, where $x_j$ is the fraction of successes over the interactions with provider $j$. The customer has a belied about each provider $j$
    \begin{align}
        \theta_A\mid x_A&\sim\operatorname{Beta}(1+\frac{x_A}{1-\rho},1+\frac{1-x_A}{1-\rho}),\\
        \theta_B\mid x_B&\sim\operatorname{Beta}(1+\frac{x_B}{1-\rho},1+\frac{1-x_B}{1-\rho})
    \end{align}
\end{frame}

\begin{frame}{Customer's choice}
    The customer uses Thompson sampling (TS), drawing independently from the
    two belief distributions and choosing
    \begin{equation}
        y_t=
        \begin{cases}
            A,&\widetilde\theta_A>\widetilde\theta_B,\\
            B,&\widetilde\theta_B>\widetilde\theta_A.
        \end{cases}
    \end{equation}
    Thus the selection probabilities are
    \begin{equation}
        \begin{aligned}
            D_A(x_A,x_B)&=\mathbb{P}(\widetilde\theta_A>\widetilde\theta_B\mid x_A,x_B),\\
            D_B(x_A,x_B)&=1-D_A(x_A,x_B).
        \end{aligned}
    \end{equation}
\end{frame}

\begin{frame}{State transitions}
    Superscripts $s$ and $f$ denote success and failure,
    respectively. The state transitions are
    \begin{equation}
        \begin{aligned}
            T_A^s(x)&=(\rho x_A+1-\rho,x_B),\\
            T_A^f(x)&=(\rho x_A,x_B),\\
            T_B^s(x)&=(x_A,\rho x_B+1-\rho),\\
            T_B^f(x)&=(x_A,\rho x_B).
        \end{aligned}
    \end{equation}
\end{frame}

\begin{frame}{Strategies}
    We consider stationary mixed strategies that depend only
    on the seller's own reputation. For each seller $j\in\{A,B\}$,
    a strategy is a mapping
    \begin{equation}
        \sigma_j:[0,1]\rightarrow\Delta(\{1,2\}),
    \end{equation}
    where $\sigma_j(i\mid x_j)$ denotes the probability that
    seller $j$ chooses product $i$ when selected at reputation
    $x_j$.
\end{frame}

\begin{frame}{Distribution-Based Oblivious Equilibrium}
    Suppose that seller $A$ assumes that its competito's reputation follows a stationary distribution $\mu_B$ on $[0,1]$. Seller $A$ computes its expected demand as 
    \begin{equation}
        d_j(x_A;\mu_{B}) =\int_0^1 D(x_A,y)\,\mu_{B}(dy)
    \end{equation}
    The seller's problem becomes 
    \begin{equation}
        \begin{aligned}
            U_A(x_A) ={}& \gamma[1-d_A(x_A;\mu_{B})]U_A(x_A)\\ &+ d_A(x_A;\mu_{B})
            \max_{i\in\{1,2\}}\Big\{ R-c_i\\
            &\qquad+\gamma\left[ p_i U_A(T_\rho^s(x_A)) + (1-p_i)U_A(T_\rho^f(x_A))
            \right] \Big\}.
        \end{aligned}
    \end{equation}
\end{frame}

\begin{frame}{Distribution-Based Oblivious Equilibrium}
    \begin{itemize}
        \item Given the conjectured distribution $\mu_B$, each seller solves its Bellman equation and obtains an optimal strategy $\sigma$.
        \item We require this distribution to be consistent with the reputations generated when both sellers follow $\sigma$.
    \end{itemize}
\end{frame}

\begin{frame}{Joint transition probability kernel}
    We define the joint transition probability kernel as
    \[
        P^\sigma(x,C)=\mathbb{P}(x_{t+1}\in C\mid x_t=x,\sigma),
    \]
    which is equivalent to
    \begin{equation}
        \begin{aligned}
            P^\sigma(x,C)
            ={}&D(x_A,x_B)\Big[
                p_{\sigma_A}(x_A)\mathbf{1}\{T_A^s(x)\in C\}\\
            &\qquad +(1-p_{\sigma_A}(x_A))\mathbf{1}\{T_A^f(x)\in C\}\Big]\\
            &+[1-D(x_A,x_B)]\Big[
                p_{\sigma_B}(x_B)\mathbf{1}\{T_B^s(x)\in C\}\\
            &\qquad +(1-p_{\sigma_B}(x_B))\mathbf{1}\{T_B^f(x)\in C\}\Big].
        \end{aligned}
    \end{equation}
\end{frame}

\begin{frame}{Distribution-Based Oblivious Equilibrium}
    A symmetric oblivious equilibrium is a triple
    $(\sigma^*,\mu^*,\nu^*)$ satisfying:

    \medskip
    \textbf{1. Optimality}
        \begin{equation*}
            \begin{aligned}
            \operatorname{supp}\sigma^*(\cdot\mid x)
            \subseteq
            \arg\max_{i\in\{1,2\}}
            \Big\{
            R-c_i+\gamma\big[
            &p_iU_{\mu^*}(T_\rho^s(x))\\
            &+(1-p_i)U_{\mu^*}(T_\rho^f(x))
            \big]
            \Big\},
            \quad \forall x\in[0,1].
            \end{aligned}
        \end{equation*}

    \medskip
    \textbf{2. Consistency}
        \begin{equation*}
            \begin{aligned}
            \nu^*&=\nu^*P^{\sigma^*},\\
            \mu^*(C)&=\nu^*(C\times[0,1])
            =\nu^*([0,1]\times C),
            \quad \forall C\subseteq[0,1]\text{ measurable}.
            \end{aligned}
        \end{equation*}
\end{frame}

\begin{frame}{Mean-Draw Oblivious Equilibrium}
    Suppose that seller $A$ represents its competitor
    by the long-run average of seller $B$'s
    Thompson sampling draws:
    \begin{equation}
        p_{B,0}
        =
        \lim_{T\to\infty}
        \frac{1}{T}
        \sum_{t=1}^{T}\widetilde{\theta}_{B,t}.
    \end{equation}

    Under stationarity and ergodicity, the time average
    coincides with the stationary expectation almost surely.
    Since
    \begin{equation}
    \mathbb{E}[\widetilde{\theta}_{B,t}\mid x_{B,t}]
    =
    \frac{1-\rho+x_{B,t}}{3-2\rho},
    \end{equation}
    we obtain
    \begin{equation}
    \begin{aligned}
    p_{B,0}
    &=\mathbb{E}_{\nu}[\widetilde{\theta}_B]\\
    &=\int_{[0,1]^2}
    \frac{1-\rho+x_B}{3-2\rho}
    \,\nu(dx_A,dx_B).
    \end{aligned}
    \end{equation}
        
\end{frame}

\begin{frame}{Mean-Draw Oblivious Equilibrium}
    Seller $A$ computes its expected demand as 
    \begin{equation}
        d_j(x_A;p_{B,0}) = \mathbb{P}(\tilde{\theta_A}>p_{B,0} \mid x_A)
    \end{equation}
    The seller's problem becomes 
    \begin{equation}
        \begin{aligned}
            U_{p_{0,j}}(x_j) ={}& \gamma[1-d(x_j;p_{0,j})]U_{p_{0,j}}(x_j)\\
            & + d(x_j;p_{0,j}) \max_{i\in\{1,2\}} \Big\{ R-c_i\\
            &\qquad+ \gamma\left[ p_iU_{p_{0,j}}(T_\rho^s(x_j)) + (1-p_i)U_{p_{0,j}}(T_\rho^f(x_j)) \right] \Big\}.
        \end{aligned}
    \end{equation}
\end{frame}

\begin{frame}{Mean-Draw Oblivious Equilibrium}

    A symmetric oblivious equilibrium is a triple
    $(\sigma^*,p_0^*,\nu^*)$ satisfying:

    \medskip
    \textbf{1. Optimality}
    \begin{equation*}
        \begin{aligned}
            \operatorname{supp}\sigma^*(\cdot\mid x)
            \subseteq
            \arg\max_{i\in\{1,2\}}
            \Big\{
            R-c_i+\gamma\big[
            &p_iU_{p_0^*}(T_\rho^s(x))\\
            &+(1-p_i)U_{p_0^*}(T_\rho^f(x))
            \big]
            \Big\},
            \quad \forall x\in[0,1].
            \end{aligned}
            \end{equation*}

            \medskip
            \textbf{2. Consistency}
            \begin{equation*}
            \begin{aligned}
            \nu^*&=\nu^*P^{\sigma^*},\\[6pt]
            p_0^*
            &=
            \int_{[0,1]^2}
            \frac{1-\rho+x_j}{3-2\rho}
            \,\nu^*(dx_A,dx_B),
            \quad j\in\{A,B\}.
        \end{aligned}
    \end{equation*}

\end{frame}

\begin{frame}{Mean-Draw: Temporal Consistency}
    The optimality condition is unchanged. For a symmetric profile
    $\sigma^*$, consistency uses the expected occupation law over a
    finite window:

    \medskip
    \begin{equation*}
        \begin{aligned}
            \bar{\nu}_{b,T}^{\,\sigma^*}
            &=\frac{1}{T}\sum_{t=b}^{b+T-1}
            \nu_0\left(P^{\sigma^*}\right)^t,\\[8pt]
            p_0^*
            &=\int_{[0,1]^2}
            \frac{1-\rho+x_j}{3-2\rho}
            \,\bar{\nu}_{b,T}^{\,\sigma^*}(dx_A,dx_B),
            \quad j\in\{A,B\}.
        \end{aligned}
    \end{equation*}

    \medskip
    \textbf{Experiment}
    \[
        \nu_0=\delta_{(0.5,\,0.5)},\qquad
        b=10\,000,\qquad T=100\,000.
    \]
    States are measured before selection. We discard the first $b$
    periods and average the next $T$ periods.

    \medskip
    The average law $\bar{\nu}_{b,T}^{\,\sigma^*}$ need not be stationary.
\end{frame}

\begin{frame}{Simulations}
    go to vscode, there are too many plots 
\end{frame}

\end{document}
